%% ***************************************************************************
%% My paper
%%
%% Authors: Emmett Brown, Marty McFly, Biff Tannen
%%
%% NOTE: this file will not compile until you called the script
%% generate-preamble.php once. See the file Readme.md to understand what
%% to do.
%%
%% This paper is an instance of the PaperShell template. For more
%% information, please visit https://github.com/sylvainhalle/PaperShell
%% ***************************************************************************
%% ---------------------------
%% Author preamble. The contents of this file change depending on the
%% paper class you choose.
%% ---------------------------
  
%%%%%%%%%%%%%%%%%%%%%%%%%%%%%%%%%%%%%%%%%%%%%%%%%%%%%%%%%%%%%%%%%%%%%%%%%%
%% This file was autogenerated by PaperShell v2.2.1 on 2020-02-26 06:57:47
%% https://github.com/sylvainhalle/PaperShell
%% DO NOT EDIT!
%%%%%%%%%%%%%%%%%%%%%%%%%%%%%%%%%%%%%%%%%%%%%%%%%%%%%%%%%%%%%%%%%%%%%%%%%%
\documentclass[preprint]{elsarticle}

% Usual packages
\usepackage[utf8]{inputenc}  % UTF-8 input encoding
\usepackage[T1]{fontenc}     % Type1 fonts
\usepackage{lmodern}         % Improved Computer Modern font
\usepackage{microtype}       % Better handling of typo
\usepackage[scaled]{helvet} % Scale Helvetica
%\usepackage[english]{babel}  % Hyphenation
\usepackage[spanish,es-noquoting]{babel}
\usepackage{graphicx}        % Import graphics
\usepackage[bookmarks=false]{hyperref}            % Better handling of references in PDFs
\usepackage{comment}         % To comment out blocks of text
\biboptions{sort&compress}   % Sort and compress citations
\usepackage[left=1in, right=1in, top=1in, bottom=1in]{geometry}
\journal{TISSEC}

% User-defined includes
%% ---------------------------
%% If you have other packages or command definitions you'd like to
%% include, write them there
%% ---------------------------
\usepackage{amsmath,amsfonts,amssymb}

%% Note: the subfig package is incompatible with the "unfixed" version
%% of the LIPICS class. It works here only because PaperShell includes
%% a fixed version --which is not the official version.
\usepackage{subfig}

%% We put one here: this is a dummy command for LaTeX, used to
%% tell Aspell to skip checking the spelling of what's inside
%% Usage: I like the word \nospellcheck{kwyjibo}
\newcommand{\nospellcheck}[1]{#1}

%% --------------------------
%% Personalized todo notes
%% --------------------------
\usepackage{xcolor}
\usepackage{todonotes}
%\newcommand{\todosylvain}[1]{\todo[inline,caption={},color=cyan]{\sf\small \textbf{@Sylvain:} #1}}
%\newcommand{\todoalex}[1]{\todo[inline,caption={},color=pink]{\sf\small \textbf{@Alex:} #1}}
%\newcommand{\todosebastien}[1]{\todo[inline,caption={},color=lime]{\sf\small \textbf{@Sébastien:} #1}}
%\newcommand{\todomichael}[1]{\todo[inline,caption={},color=lime]{\sf\small \textbf{@Michaël:} #1}}
%\newcommand{\todotous}[1]{\todo[inline,caption={},color=yellow]{\sf\small #1}}
%\newcommand{\todoedmond}[1]{\todo[inline,caption={},color=lime]{\sf\small \textbf{@Edmond:} #1}}

%% --------------------------
%% Placeholder for figure
%% --------------------------
\newcommand{\imagevide}{\framebox(200,200){IMAGE}}

%% --------------------------
%% Alter some LaTeX defaults for better treatment of figures:
%% See p.105 of "TeX Unbound" for suggested values.
%% See pp. 199-200 of Lamport's "LaTeX" book for details.
%% --------------------------
%% Parameters for all pages
\renewcommand{\topfraction}{0.9}	% max fraction of floats at top
\renewcommand{\bottomfraction}{0.8}	% max fraction of floats at bottom

%% Parameters for TEXT pages (not float pages)
\setcounter{topnumber}{2}
\setcounter{bottomnumber}{2}
\setcounter{totalnumber}{4}                 % 2 may work better
\setcounter{dbltopnumber}{2}                % for 2-column pages
\renewcommand{\dbltopfraction}{0.9}         % fit big float above 2-col. text
\renewcommand{\textfraction}{0.07}          % allow minimal text w. figs

% Parameters for FLOAT pages (not text pages)
\renewcommand{\floatpagefraction}{0.7}      % require fuller float pages
% N.B.: floatpagefraction MUST be less than topfraction !!
\renewcommand{\dblfloatpagefraction}{0.7}   %require fuller float pages

%\newtheorem{example}{Example}
%% Default path for graphics
\graphicspath{{fig/}}

\date{Julio, 2020}
\begin{document}

\begin{frontmatter}

% Title
%\title{PRINCIPALES FACTORES DE VIOLENCIA CONTRA LAS MUJERES}
\title{{\Large \textbf{PRINCIPALES FACTORES DE VIOLENCIA CONTRA LAS MUJERES}}}
% Authors and affiliations
\author{ Enrique Alejandro Laurel Cossío\\--\\Docente: Alvaro Limber Chirino Gutierrez\\Materia: Programación II\\ Universidad Mayor de San Andrés}
%\affil{Docente: Alvaro Limber Chirino Gutierrez\\Materia: Programación II\\ Universidad Mayor de San Andrés}
%\author{Emmett Brown\fnref{label1}%
%}
%\author{Marty McFly\fnref{label1}%
%}
%\author{Biff Tannen\fnref{label2}%
%}
\fntext[label1]{Temporal Industries, Hill Valley, CA 95420%
}
\fntext[label2]{BiffCo inc., Hill Valley, CA 95420%
}
%\begin{abstract}
%%% ----------------------
%% Write your abstract here. Do not enclose it in an "abstract"
%% environment.
%% ----------------------
Lorem ipsum dolor sit amet, consectetur adipiscing elit. Quisque rhoncus auctor mauris. Vestibulum ante ipsum primis in faucibus orci luctus et ultrices posuere cubilia Cur\ae{}; Sed a metus nec felis tristique venenatis eu quis urna. Integer quis sagittis eros, condimentum auctor libero. Fusce luctus diam et libero sodales luctus. Duis suscipit sodales libero, sed blandit odio aliquam at. Sed venenatis orci et luctus sagittis. Nam aliquam mi mi, ut vehicula turpis semper at. Nam pellentesque est nec fringilla consectetur. Aliquam magna elit, imperdiet et condimentum nec, venenatis et nisi. 
%% :folding=explicit:wrap=soft:mode=latex:

%\end{abstract}
\end{frontmatter}


%% ---------------------------
%% If you wish to include additional packages, define new environments or
%% new commands, put them in the file includes.tex
%%
%% Write your abstract in the file abstract.tex.
%% ---------------------------

%% ---------------------------
%% Introduction
%% ---------------------------

\section{Introducción} %% {{{

La violencia contra las mujeres es uno de los principales problemas sociales de nuestro país, que a lo largo de nuestra historia republicana ha sido relegada al ámbito privado, hasta la aprobación de nuestra Constitución Política del Estado el 2009, donde la Ley Fundamental
establece como valor del Estado la equidad de género, y prohíbe cualquier tipo de discriminación en razón de género.

Violencia contra la mujer es la que se ejerce por su condición de mujer. Siendo esta «consecuencia de la discriminación que sufre tanto en leyes como en la práctica, y la persistencia de desigualdades por razones de género».

En esta violencia se presentan numerosas facetas que van desde la discriminación y el menosprecio hasta la agresión física, sexual, verbal, psicológica y el asesinato, manifestándose en diversos ámbitos de la vida social, laboral y política, entre los que se encuentran la propia familia, la escuela, las religiones, el Estado, entre otras.

En 1993, en asamblea general, las Naciones Unidas (ONU) aprobaron la Declaración sobre la eliminación de la violencia contra la mujer, y en 1999, a propuesta de la República Dominicana con el apoyo de 60 países más, declararon el 25 de noviembre Día Internacional de la Eliminación de la Violencia contra la Mujer.

En 2008 el Secretario General de la ONU puso en marcha la campaña «Unidos para poner fin a la violencia contra las mujeres» apelando al «imperio de la ley» como vehículo para su erradicación. Uno de sus objetivos fue el de procurar que para 2015 todos los países hubiesen adoptado leyes específicas contra este tipo de violencia de conformidad con las normas internacionales en materia de derechos humanos.

En febrero de 2008 el Secretario General de Naciones Unidas Ban Ki-moon lanzó la campaña ÚNETE para poner fin a la violencia contra las mujeres, proclamando el 25 de cada mes Día Naranja. Entre otras actividades, en ese día se invita a llevar alguna prenda de ese color para resaltar el llamado a erradicar la violencia contra la mujer.

El trabajo fue realizado con la base de datos proporcionados por el Instituto Nacional de Estadística, primera ?Encuesta de Prevalencia y Características de la Violencia contra las Mujeres? (EPCVcM).

Los avances en los últimos años son significativos, en cuanto a Tratados y Convenios Internacionales sobre Derechos Humanos suscritos y ratificados por el Estado Plurinacional de Bolivia; la Ley Integral para Garantizar a las Mujeres una Vida Libre de Violencia N° 348


\section{Objetivos}

\subsection{Objetivo General}
Encontrar los principales factores o razones de la violencia contra las mujeres
\subsection{Objetivos Específicos}
\begin{itemize}
\item Indicios o indicadores de feminicidios en Bolivia 
\end{itemize}

\section{Motivación}
Se menciona en noticias de ámbito nacional que los casos reportados de violencia contra la mujer y situaciones peores como feminicidios se elevaron en esta cuarentena por razones del COVID-19. 
Razón por la cual quiero mostrar los factores mas importantes de esta problemática.
Factores que puedan ayudar a tomar decisiones o nuevas políticas.

\section{Marco Teórico}

\subsection{Marco Normativo}
La Constitución Política del Estado (CPE) es nuestra Ley Fundamental mediante la cual se refunda el Estado con la voluntad de todas y todos los ciudadanos, ha replanteado las estructuras jurídicas, económicas y sociales de Bolivia, así que la misma CPE, establece mandatos a ley
específicos, debiendo además adecuar toda la legislación
anterior a la CPE.

En la CPE la Equidad de Género se desarrolla transversalmente en todo el texto constitucional:
\begin{itemize}
\item Como valor del Estado, Equidad de género en la participación (Art 8. II)
\item Sistema de Gobierno con equivalencia entre hombres y mujeres (Art. 11. I)
\item Prohibición y sanción a toda forma de discriminación (Art.14. II)
\item Derecho Fundamental a la Vida, Prevención y Sanción de la Violencia de Género (Art. 15 II y III)
\item Derechos Políticos, en igualdad de condiciones entre hombres y mujeres (Art. 26)
\item Representación Política (Art. 210. II)
\item Participación Política- designación de ministros (Art. 172 - 22)
\item Elección de autoridades Departamentales (Art. 278) Asambleístas nacionales (Art. 147)
\item Educación- Valores de equidad de Género (Art. 79)
\item Principio de descentralización y organización territorial (Art. 270)
\end{itemize}
 
Dentro de la normativa jurídica establecida por la CPE, en cumplimiento de los mandatos de las distintas Convenciones y Tratados que Bolivia firmó, existen leyes y decretos promulgados durante los últimos 5 años relacionados a la protección de los derechos de las mujeres, en especial, al de una Vida Libre de Violencia, entre las más importantes:

Ley Integral para Garantizar a las Mujeres una Vida Libre de Violencia, Ley N° 348 de 9 de marzo de 2013, en el Artículo 1 establece ?(MARCO CONSTITUCIONAL). La presente Ley se funda en el mandato constitucional y en los Instrumentos, Tratados y Convenios Internacionales de Derechos Humanos ratificados por Bolivia, que garantizan a todas las personas, en particular a las mujeres, el derecho a no sufrir violencia física, sexual y/o psicológica, tanto en la familia como en la sociedad.?
Tiene por objeto establecer mecanismos, medidas y políticas integrales de prevención, atención, protección y reparación a las mujeres en situación de violencia, así como, la persecución y sanción a los agresores, con el fin de Garantizar a las Mujeres una Vida Digna y el ejercicio pleno de sus derechos para Vivir Bien.
Esta Ley, además establece en el Artículo 11 la necesidad de contar con información en el Sistema Integral Plurinacional de Prevención, Atención, Sanción y Erradicación de la Violencia en Razón de Género ? SIPPASE-VRG.
El SIPPASE tiene como mandato construir e Implementar todo el Sistema de Información que generan las instituciones públicas y privadas como fuente de información estadística sobre esta problemática, tarea que está siendo diseñada entre el Ente Rector, el Viceministerio de Igualdad de Oportunidades, a través de la Dirección General de Prevención y Erradicación de Toda Forma de Violencia en Razón de Género y Generacional, del Ministerio de Justicia y Transparencia Institucional y el
Instituto Nacional de Estadística.

\subsection{Definiciones y Conceptos}
La violencia contra las mujeres, es y será un importante campo de estudio, que incluye investigaciones técnicas confiables y válidas para medir esta problemática social, que en gran parte de la región ya está incorporada en sus normativas jurídicas y, varían de un país a otro.
Unas de las primeras definiciones que se dieron sobre violencia contra las mujeres, fue la enunciada por las Naciones Unidas, en el marco de la Asamblea General de 1993:
Violencia contra la mujer: "Todo acto de Violencia basado en la pertenencia al sexo femenino que tenga o pueda tener como resultado un daño o sufrimiento físico, sexual o psicológico para la mujer, así como las amenazas de tales actos, la coacción o la privación arbitraria de la libertad, tanto si se producen en la vida pública como en la vida privada ''
Así también se tiene la definición de la "Convención de Belém do Pará'', mencionada en Capítulos anteriores, que señala que debe entenderse por violencia contra la Mujer "cualquier acción o conducta, basada en su género,
que cause muerte, daño o sufrimiento físico, sexual o psicológico a la mujer, tanto en ámbito público como en el privado''.
En ambas definiciones se identifica una variedad de modalidades bajo las cuales se presenta la violencia contra las mujeres y se diferencian los ámbitos en que ésta se desarrolla; sin embargo, para poder medirla con todas sus particularidades, recurriremos a las definiciones textuales que da la Ley N° 348, que conceptualiza los distintos tipos y formas de violencia.
"Violencia. Constituye cualquier acción u omisión, abierta o encubierta, que cause la muerte, sufrimiento o daño físico, sexual o psicológico a una mujer u otra persona, le genere perjuicio en su patrimonio, en su economía, en su fuente laboral o en otro ámbito cualquiera, por el sólo hecho de ser mujer''. (Art. 6.1)
"Situación de violencia. Es el conjunto de circunstancias y condiciones de Agresión en las que se encuentra una mujer, en un momento determinado de su vida''. (Art.6.2)
En la presente investigación se ha tomado en cuenta cuatro tipos de violencia contra las mujeres: física, psicológica, sexual, económica y/o patrimonial, los cuales dependiendo del ámbito donde se manifiesten, sea este privado o público, pueden asumir diferentes formas. En este sentido, empecemos por definir qué se entiende por ámbito privado y ámbito público:

\subsubsection{Ámbito Privado}
Está referido al espacio donde se desenvuelve la familia, la unidad doméstica y también, cualquier otra relación interpersonal.
La violencia contra las mujeres se da con más intensidad en este ámbito por ser el lugar más común para su ejecución. Se expresa en los cuatro tipos de violencia señalados anteriormente, que se establece en la Ley 348:
Violencia Física. Es toda acción que ocasiona lesiones y/o daño corporal, interno, externo o ambos, temporal o permanente, que se manifiesta de forma inmediata o en el largo plazo, empleando o no fuerza física, armas o
cualquier otro medio.? (Art. 7.1)
Violencia Sexual. Es toda conducta que ponga en riesgo la autodeterminación sexual, tanto en el acto sexual como en toda forma de contacto o acceso carnal, genital o no genital, que amenace, vulnere o restrinja el derecho al ejercicio a una vida sexual libre segura, efectiva y plena, con autonomía y libertad sexual de la Mujer. (Art. 7.7)
Violencia Psicológica. Es el conjunto de acciones sistemáticas de desvalorización, intimidación y control del comportamiento, y decisiones de las mujeres, que tienen como consecuencia la disminución de su autoestima, depresión, inestabilidad psicológica, desorientación e incluso el suicidio. (Art. 7.3) 
Estos cuatro Tipos de Violencia no se los debe ver de forma aislada sino de manera combinada, pues a través de diferentes investigaciones se ha demostrado que en los casos de Violencia Extrema que sufren algunas mujeres a lo largo de sus vidas, han estado presentes todos estos tipos, y muchas de ellas se encontraban en una situación de alto riesgo. 

\subsubsection{Ámbito Público}
Todos los tipos de violencia mencionados, también se dan en el ámbito público, el cual se entiende como el espacio externo a las fronteras de la familia, donde se establecen las relaciones sociales que no están mediadas
por el parentesco, como son: las unidades educativas, el trabajo, la calle, los lugares de esparcimiento social, entre otros, donde se patentiza el poder del género masculino sobre el femenino como una característica del sistema
patriarcal y una muestra de la violencia estructural contra las mujeres.
Para el presente caso de estudio se identifica: el espacio educativo; el espacio laboral; los servicios de salud; la calle, que es el espacio de la cotidianidad pública. 
La violencia contra las mujeres en el espacio educativo, suele expresarse de distintas formas: discriminación, acoso verbal y sexual intimidando a la víctima o castigándola por medio de agresiones (caricias no deseadas, relaciones sexuales forzadas, condicionamientos, etc.) Establecidos
en la Ley 348:
?Violencia en el Sistema Educativo Plurinacional. Es todo acto de agresión física, psicológica o sexual cometido contra las mujeres en el sistema educativo regular, alternativo, especial y superior.? (Art. 7.12)
En el espacio laboral las manifestaciones de la violencia hacia las mujeres se presentan como: acoso, hostigamiento sexual, segregación, discriminación salarial, mayores restricciones de contratación por su estado civil, gravidez, etc. y relegación a tareas subordinadas y de servicio, entre
otras. "Violencia Laboral. Es toda acción que se produce en cualquier ámbito de trabajo por parte de cualquier persona de superior, igual o inferior jerarquía que discrimina, humilla, amenaza o intimida a las mujeres;
que obstaculiza o supedita su acceso al empleo, permanencia o ascenso y que vulnera el ejercicio de sus derechos.'' (Art. 11)
Al constituirse la violencia hacia las mujeres en un problema de salud pública se han dado competencias específicas a este sector, tanto en lo que se refiere a educación y promoción como a atención. En el espacio de la salud se ve a diario violencia contra las mujeres, a quienes muchas veces se da un trato deshumanizado que no respeta su intimidad e incluso se discrimina a mujeres por ser pobres o indígenas.
Violencia en Servicios de Salud. Es toda acción discriminadora, humillante y deshumanizada y que omite, niega o restringe el acceso a la atención eficaz e inmediata y a la información oportuna por parte del personal de salud, poniendo en riesgo la vida y la salud de las mujeres.? (Art. 7.9)
El derecho de las mujeres a tener información sobre su salud sexual y salud reproductiva es fundamental para prevenir embarazos no deseados y enfermedades e infecciones de transmisión sexual. La violencia en servicios
de Salud, es parte de la lucha patriarcal por dominar el cuerpo de las mujeres, al no permitir que éstas gocen de su sexualidad o de su vida reproductiva.
Violencia Contra los Derechos Reproductivos. Es la acción u omisión que impide, limita o vulnera el derecho de las mujeres a la información, orientación, atención integral y tratamiento durante el embarazo o pérdida, parto, puerperio y lactancia; a decidir libre y responsablemente el número y espaciamiento de hijas e hijos; a ejercer su maternidad segura, y a elegir métodos anticonceptivos seguros.? (Art. 7.8)
El espacio social es también donde se manifiesta la violencia contra las mujeres de muchas maneras; éste es el lugar de la cotidianidad, como la calle, el cine, las fiestas, etc. Aquí es donde se da una violencia con expresiones variadas como insultos, piropos obscenos, manoseos impúdicos e intimidaciones de tipo sexual; las formas más dolorosas de la violencia callejera son la violación sexual y el feminicidio.

\subsection{Software}
R software es un entorno y lenguaje de programación diseñado para el análisis estadístico.

Puedes: estudiar correlaciones, ajustar modelos, crear gráficos 3D de altísima calidad, aplicar árboles de decisión, realizar análisis clúster, análisis de componentes principales, crear redes neuronales de predicción, etc.
La lista de posibilidades es muy pero que muy extensa y se adapta a todo tipo de necesidades para el análisis complejo de datos.

\section{Descripción de la Base de Datos}
\subsection{Origen}
El presente documento hace referencia a la primera "Encuesta de Prevalencia y Características de la Violencia contra las Mujeres'' (EPCVcM).
La encuesta y/o recolección de la información se la realizo en cuatro cuestionarios, los cuales son:
\begin{itemize}
\item	Cuestionario Violencia Hogar.
\item   Cuestionario Violencia Casada.
\item   Cuestionario Violencia Soltera.
\item   Cuestionario Violencia Separada.
\end{itemize}
Se cuenta con 5 bases de datos:
\begin{itemize}
\item	Casadas
\item	Solteras
\item	Persona
\item	Separadas
\item	Vivienda
\end{itemize}

\subsection{Variables de Interés}
\begin{itemize}
\item ¿ Cuántos años cumplidos tiene?
\item ¿ Cuál es su estado civil o conyugal actual?
\item Departamento
\item Área geográfica
\item ¿ A lo largo de su vida, en algún lugar público algún(os) hombre(s) conocido(s) o desconocido(s), sin considerar a su (ex) pareja/ (ex) novio, la han humillado o menospreciado (la han hecho sentir menos)?
\item variables de la base de datos mujeres casadas, sección 3, pregunta 10
\item ¿ Usted estudia o ha estudiado en un centro educativo (escuela, colegio, universidad u otro)?
\item ¿ Usted trabaja o trabajo para alguien de su familia, para otra persona o por cuenta propia?
\item variables de la base de datos mujeres separadas, sección 2, pregunta 5
\item variables de la base de datos mujeres separadas, sección 3, pregunta 10
\end{itemize}

%% }}} --- Section

%% ---------------------------
%% A section
%% ---------------------------
\section{Donec id eros non nisl pharetra} %% {{{

In hendrerit commodo urna sit amet egestas. Phasellus ut faucibus diam. Etiam quis hendrerit augue. Nam vel arcu at massa iaculis ullamcorper. Sed in malesuada enim, ac rutrum augue \cite{Lenat83a}. Donec efficitur egestas massa in varius. Etiam at nibh commodo, iaculis massa nec, rutrum sem. Mauris imperdiet massa eu nibh fermentum, ac ultricies metus sollicitudin. Donec vel dolor non turpis efficitur rhoncus. M\ae{}cenas et augue congue elit viverra sagittis quis in magna. Vestibulum tempus tellus in efficitur viverra. Morbi vestibulum posuere tortor, vit\ae{} eleifend dolor iaculis id. Vestibulum ornare gravida tortor vel fermentum. Cras vulputate facilisis dui, non porttitor arcu blandit vel.

\begin{figure}
  \centering
  \includegraphics[width=0.4\textwidth]{square-circle}
  \caption{A square and a circle}
  \label{fig:square-circle}
\end{figure}

In maximus ante at metus vulputate congue. Sed eleifend ultrices tincidunt. Lorem ipsum dolor sit amet, consectetur adipiscing elit. Pr\ae{}sent quis rutrum elit. Nunc auctor nibh non sem consequat gravida. Nulla non dapibus metus. Cras viverra tempor nibh sit amet malesuada. Etiam non ante purus. M\ae{}cenas eleifend ultricies orci, ut fermentum erat suscipit elementum. Sed est libero, laoreet vel sodales ac, consectetur sed enim. Quisque hendrerit ac erat vit\ae{} posuere. Sed eros sapien, luctus ut justo lacinia, vestibulum sagittis lorem. Mauris porta dapibus dui, eu iaculis dolor faucibus id.

%% }}} --- Section

%% ---------------------------
%% Bibliography and postamble
%% ---------------------------
%%%%%%%%%%%%%%%%%%%%%%%%%%%%%%%%%%%%%%%%%%%%%%%%%%%%%%%%%%%%%%%%%%%%%%%%%%
%% This file was autogenerated by PaperShell v2.2.1 on 2020-02-26 06:57:47
%% https://github.com/sylvainhalle/PaperShell
%% DO NOT EDIT!
%%%%%%%%%%%%%%%%%%%%%%%%%%%%%%%%%%%%%%%%%%%%%%%%%%%%%%%%%%%%%%%%%%%%%%%%%%
\bibliographystyle{elsarticle-num}

\section*{References}

\bibliography{paper}
%%%%%%%%%%%%%%%%%%%%%%%%%%%%%%%%%%%%%%%%%%%%%%%%%%%%%%%%%%%%%%%%%%%%%%%%%%
%% This file was autogenerated by PaperShell v2.2.1 on 2020-02-26 06:57:47
%% https://github.com/sylvainhalle/PaperShell
%% DO NOT EDIT!
%%%%%%%%%%%%%%%%%%%%%%%%%%%%%%%%%%%%%%%%%%%%%%%%%%%%%%%%%%%%%%%%%%%%%%%%%%
%% ---------------------------
%% If there is anything you would like to insert *after* the bibliography
%% (and *before* the \end{document} instruction), place it here
%% ---------------------------

% Nothing



\end{document}

%% :folding=explicit:wrap=soft:mode=latex: